\chapter{Neun Zählfolgen aus der Horimetrik}
\label{ch:zaehlungen}

% Register: D8, D9, S9, S19, S22, S41, V11. Quellen: Anhang A.4-A.6,
% C.6, F.1-F.5, H, oeis/ENTWUERFE.md.

Theorien hinterlassen Fingerabdrücke: Zahlenfolgen, die beim Zählen
ihrer Objekte entstehen. Die Horimetrik hat im Lauf dieses Projekts
neun solcher Folgen erzeugt, und keine davon findet sich in der
OEIS, der großen Enzyklopädie der Zahlenfolgen -- was kein Beweis
für Neuheit ist, aber ein Indiz, dass hier bisher niemand gezählt
hat. Alle neun werden bei der OEIS eingereicht; die zugeteilten
A-Nummern tragen wir nach Vergabe hier nach.\footnote{Platzhalter:
A??????. Die Einreichungsentwürfe mit jargonfreien Definitionen und
Programmen liegen im Projektordner \texttt{oeis/}.}

\section{Die Steckbriefe}

\paragraph{Die Sprungstellen des Tiefengesetzes.}
$3,\ 15,\ 84,\ 533,\ 3883,\ 31998,\ 294875,\ 3006206,\ \ldots$
Die Schale, in der die Kommutator-Tiefe
(Satz~\ref{satz:tiefengesetz}) erstmals $-g$ erreicht.
Die Königin unserer Folgen: Für sie ist ein asymptotisches Gesetz
\emph{bewiesen} ($m_g/(g+2)!\to1$), und sie hat drei Vorhersagetests
bestanden -- den letzten exakt: $m_7=294875$ wurde prognostiziert,
bevor es berechnet war.

\paragraph{Die maximale Produkthöhe von Quadraten.}
$1,\ 6,\ 22,\ 87,\ 407,\ 2473,\ 19527,\ \ldots$
Wie groß ein Horizont werden kann, wenn man zwei randvolle
Gestalten der Kapazität $r$ multipliziert (Register D8). Inzwischen bewiesen: Die Folge wächst wie
$r!\cdot e^{2\sqrt r+O(\log r)}$ -- eine Fakultät mal einem sanft
explodierenden Faktor (Anhang~\ref{ch:anhang-beweise}); nur die
Feinstruktur des Logarithmus-Terms ist noch offen.

\paragraph{Die Schalenspringer.}
$1,\ 8,\ 60,\ 360,\ 2520,\ 21168,\ 211680,\ 2268000,\ \ldots$
Die Anzahl der Zahlen, die eine einzige führende Null in die
nächste Fakultätsschale hebt. Diese Folge hat uns das
$n!/2$-Debakel beschert (Kapitel~\ref{ch:sackgassen}): Drei Terme
lang sah sie aus wie eine Halbfakultät, beim vierten nicht mehr.
Ihr Komplement, die \emph{Nicht-Springer} $n!-D(n)$
($5,\ 16,\ 60,\ 360,\ 2520,\ 19152,\ \ldots$), zählt als eigene
Folge der Sammlung -- ausgezeichnet durch eine vollständig
bewiesene geschlossene Formel, die Ziffernformel aus
Anhang~\ref{ch:anhang-beweise}.

\paragraph{Die beidseitigen Fixpunkte.}
$4,\ 17,\ 88,\ 540,\ 3960,\ 32760,\ \ldots$
Zahlen, deren Darstellung unter beiden Aufräum-Operationen
gleichzeitig stabil ist (Abschnitt~\ref{sec:kommutator}); zählbar
als $n\cdot n!$ minus Schalenspringer, denn von den $n\cdot n!$
wertkompakten Elementen einer Schale scheitern an der
Strukturkompaktheit genau die ohne Anschlag-Ziffer -- und das sind
die Springer-Folgen.

\paragraph{Die Ein-Ziffern-Fakultäten.}
$1,\ 2,\ 5,\ 7,\ 11,\ 23,\ 29,\ 39,\ 59,\ 119,\ 143,\ \ldots$
Die $n$, für die $n!$ in seiner eigenen Schale eine einzelne Ziffer
ist -- vollständig charakterisiert als $n=(i+1)!/d-1$ mit
$1\le d\le i$, genau $i$ Stück pro Position
(Anhang~\ref{ch:anhang-beweise}); die Position $i=1$ liefert den
trivialen Anfangsterm $n=1$. Die einzige unserer Folgen, die
schon bei ihrer Entdeckung eine geschlossene Beschreibung mitbrachte
-- und der Lieferant von neun exakt eingetroffenen Vorhersagen.

\paragraph{Die Kompositions-Explosion.}
$1,\ 4,\ 23,\ 6541,\ 477\,149\,751,\ 26\,594\,371\,862\,819\,905,\ \ldots$
Wie hoch die Kapazität klettern kann, wenn man eine randvolle
Gestalt der Kapazität $r$ in eine ebensolche \emph{einsetzt}
(Kapitel~\ref{ch:strukturpolynome}). Die jüngste und mit weitem
Abstand am schnellsten wachsende Folge der Sammlung: Einsetzen
multipliziert die Grade und lässt die Koeffizienten explodieren.
Dank Extremalprinzip trotzdem mühelos berechenbar -- ein einziges
kanonisches Einsetzen pro Term. Ihre Asymptotik ist offen.

\paragraph{Die Interchange-Maxima und -Nullzähler.}
$0,\ 1,\ 6,\ 43,\ 351,\ 2873,\ 26443,\ 270291,\ \ldots$ und
$2,\ 5,\ 13,\ 23,\ 61,\ 111,\ 274,\ 564,\ \ldots$
Die beiden Wachstumsoperationen -- erst Null voran, dann werttreu
umziehen, oder umgekehrt -- vertauschen nicht. Den Unterschied
misst der \emph{Interchange-Defekt} $\Omega$: der Wert nach „erst
wachsen, dann liften“ minus der Wert nach „erst liften, dann
wachsen“. Dieser Defekt ist \emph{niemals negativ}: Erst wachsen,
dann liften liefert nie weniger. Das kleinste
lehrreiche Beispiel ist die Drei als $10$ in $H_2$: Erst wachsen:
$010$ in $H_3$ ist die \emph{Vier}; werttreu geliftet bleibt es die
Vier. Erst liften: die Drei wird zu $003$; eine Null davor ist
gratis, es bleibt die Drei. Also $\Omega=4-3=1$. Das war lange Vermutung und ist inzwischen
Satz -- der Beweis läuft über ein Multiplikationslemma mit scharfer
Regime-Grenze. Auch die Gleichheitsfälle sind inzwischen
charakterisiert: gleichgültig gegenüber der Reihenfolge ist eine
Zahl genau dann, wenn beide Beweisschritte scharf sind --
übertragsfrei plus scharfes Lemma -- und der Nullzähler gehorcht
einer daraus abgeleiteten semi-geschlossenen Zählformel (exakt
bestätigt für alle gemessenen Terme).

\section{Warum Folgen einreichen?}

Aus zwei Gründen, einem bescheidenen und einem strategischen. Der
bescheidene: Es ist ein kleiner, echter Dienst an der Mathematik --
saubere Definitionen, Programme, ein bewiesenes Gesetz. Der
strategische: Die OEIS ist der beste Neuheitsdetektor der Welt.
Wer irgendwo, in welchem Kontext auch immer, künftig
$3, 15, 84, 533$ ausrechnet, findet diese Einträge -- und damit
dieses Buch. Umgekehrt verlinken die Einträge hierher. Sollte eine
der Folgen doch eines Tages als bekannte Größe in Verkleidung
enttarnt werden, ist auch das ein Ergebnis: Dann wissen wir endlich
präzise, welcher Teil der Horimetrik alte Mathematik in neuen
Kleidern ist -- und welcher nicht.

\begin{oberflaeche}
Neun Zählfragen, neun Antwortreihen, die offenbar noch niemand
aufgeschrieben hat. Vielleicht, weil sie neu sind. Vielleicht nur,
weil nie jemand gefragt hat. Der Unterschied zwischen diesen beiden
Möglichkeiten ist genau das, was eine Einreichung klärt -- und
darum reichen wir ein.
\end{oberflaeche}
